\begin{center}
Problema C

Abobonagramator

{\tiny Por Rafael Castro}

Timelimit: $1$	
\end{center}			
			
Após anos de intenso desenvolvimento, os cientistas do instituto 
\textit{edrev oir gfi}
finalmente lançaram sua invenção que promete ser revolucionária: O Abobonagramator. Essa máquina é capaz de transformar qualquer palavra com pelo menos 7 letras, em um anagrama de "abobora". Uma palavra é anagrama de outra quando podemos reordenar as letras da primeira até virar a segunda, por exemplo "edrev" e "verde" são anagramas, assim como "oir" e "rio".
 
    A máquina funciona realizando duas operações básicas, sobre qualquer palavra $T$, sucessivamente:

    \begin{enumerate}
        \item Substituição de uma letra por outra; ou seja, escolhendo uma letra $c$ e um índice $i$ e atribuindo o valor $T_i = c$. Por exemplo, se $T = "\textbf{a}bacate"$ e for escolhida a letra em negrito para ser substituída por 'e', então a nova palavra será $"\textbf{e}bacate"$. Essa operação custa 5 unidades de bateria da máquina. 
        
        \item Exclusão de uma letra; ou seja, escolhendo um índice $i$, e excluindo a letra $T_i$, deslocando as outras letras a frente de $i$ uma posição. Por exemplo, se $T = "u\textbf{v}auvauva"$ e for escolhida a letra em negrito para ser excluída, então a nova palavra será $"uauvauva"$. Essa operação custa 3 unidades de bateria da máquina.
    \end{enumerate}

    Dessa forma, calcule qual a mínima quantidade de unidades de bateria que o Abobonagramator precisa gastar para transformar uma palavra $S$ qualquer em um anagrama da palavra "abobora".

\vspace{0.3cm}
\textbf{Entrada}

A entrada consistirá de uma única linha contendo uma palavra $S$, contendo letras do alfabeto (a..z), minúsculas e sem acento.
\begin{itemize}
\item $7 \leqslant |S| \leqslant 10^{5}$;
\end{itemize}

\vspace{0.3cm}
\textbf{Saída}

A saída deve conter um único inteiro: A menor quantidade de unidades de bateria que a máquina gasta para transforma $S$ em um anagrama de "abobora". É garantido que sempre existe uma forma de transformar a palavra em um anagrama de "abobora".

\begin{table}[h]
    \centering
    \begin{tabular}{|p{7cm}|p{7cm}|}
        \hline
        \begin{center}
            \vspace{-0.3cm}
            \textbf{Exemplos de Entrada}
            \vspace{-0.3cm}
        \end{center} 
         &  
        \begin{center}
            \vspace{-0.3cm}
            \textbf{Exemplos de Saída}
            \vspace{-0.3cm}
        \end{center} \\ \hline
        \texttt{abbobora} & \texttt{3} \\ \hline
        \texttt{baoborex} & \texttt{8} \\ \hline
        \texttt{uvauvauva} & \texttt{28} \\ \hline
    \end{tabular}
    \caption{Exemplos de entradas e saídas}
    \label{tab:inputc2}
\end{table}

\newpage

\begin{center}
\noindent
\lstinputlisting{abobora_23/C/C3.cpp}
\end{center}